Este problema adquiere especial nitidez para los acoplamientos adimensionales.
Una coordenada expresada en unidades depende de la convención metrológica;
una relación adimensional permite preguntar por su determinación con
independencia de esa elección. La constante de estructura fina interviene
en relaciones cuantitativas contrastables y suscita, a la vez, la pregunta
por la estructura que fija el acoplamiento electromagnético. El programa
histórico de su explicación proporciona un antecedente preciso para
estudiar conjuntamente el valor de una constante y la organización
matemática de la que procede.

Un episodio esclarecedor aparece en el trabajo de Dirac de 1931 sobre
singularidades cuantizadas del campo electromagnético. Su propósito se
dirige expresamente a la existencia de una carga eléctrica mínima. El
resultado establece una relación entre esa carga y la intensidad de un
monopolo magnético; Dirac advierte, sin embargo, que la construcción no
determina por sí sola el cociente adimensional vinculado al valor
aproximado \(137\)\cite{dirac1931}. Se obtiene una restricción estructural
genuina cuya fuerza demostrativa no equivale a la fijación independiente
del acoplamiento. Este precedente muestra que la pregunta por el origen
de las constantes formaba parte de la investigación fundamental y exige
identificar exactamente qué determina cada teoría.

La historia de la electrodinámica cuántica precisa la diferencia entre
esa demanda y el cálculo de efectos observables. En su conferencia Nobel
de 1965, Feynman reconstruyó procedimientos que expresaban desplazamientos
de niveles y procesos de dispersión mediante la masa experimental del
electrón\cite{feynman1965}. La renormalización hizo posible obtener
resultados finitos y contrastables conservando esa referencia. La eficacia
de estas predicciones constituye un logro teórico efectivo; su dependencia
de un parámetro experimental identifica, a la vez, el lugar lógico en el
que una determinación adicional podría ampliar la explicación. La
cuestión constitutiva se añade a la capacidad predictiva y estudia las
condiciones que fijan sus parámetros.

El informe conjunto LEP--SLD de 2006 ofrece una manifestación experimental
de esta arquitectura\cite{lep2006}. Sus autores distinguen los parámetros
libres del modelo de las relaciones que éste impone entre observables.
Mediante correcciones radiativas, los datos de la resonancia Z permitían
inferir indirectamente las masas del quark top y del bosón W y contrastarlas
con mediciones directas. Las predicciones entrelazadas y la determinación
experimental de parámetros formaban un mismo procedimiento de prueba.
El problema relevante es, por tanto, identificar las dependencias que una
construcción determina internamente y aquellas cuya especificación sigue
requiriendo información empírica.

En este contexto, se entenderá por \emph{determinación epistemológica}
el procedimiento que establece un valor mediante la relación entre teoría,
observación e inferencia; por \emph{determinación constitutiva}, la
deducción de ese valor a partir de las relaciones que definen el objeto
considerado. Esta segunda expresión designa aquí una exigencia matemática:
explicitar el dominio, los operadores, las normalizaciones y las
compatibilidades que seleccionan el resultado. Su relevancia física
procede de la confrontación posterior con las magnitudes observadas.
El cambio de perspectiva afecta al orden de la explicación: el valor se
examina como consecuencia de una estructura determinada. La aspiración
ontológica consiste en relacionar esa estructura con la constitución del
objeto físico; la prueba matemática y su contraste experimental establecen
el alcance con el que puede sostenerse tal relación.

